\documentclass[UTF8,12pt]{ctexart}
\usepackage[a4paper,margin=24mm]{geometry}
\usepackage{amsmath,amssymb,bm,graphicx,adjustbox}
\newcommand{\fitmath}[1]{\begin{adjustbox}{max width=\linewidth}$\displaystyle #1$\end{adjustbox}}
\setlength{\parindent}{0pt}
\setlength{\parskip}{.7em}
\sloppy
\allowdisplaybreaks
\begin{document}
\section*{2017 年考研数学三 · 参考答案}
\subsection*{原 PDF 第 15 页}

\subsection*{2017年真题参考答案}

\subsection*{一、选择题}

（1）A。 （2）D。 （3）C。 （4）C。 （5）A。 （6）B。 （7）C。 （8）B。


\subsection*{二、填空题}

（9）\(\displaystyle \dfrac{\displaystyle \pi^3}{\displaystyle 2}\)。 （10）\(\displaystyle (C+\dfrac{\displaystyle 1}{\displaystyle 2}t)2^t\)。 （11）\(\displaystyle 1+e^{-Q}-Qe^{-Q}\)。 （12）\(\displaystyle xye^y\)。 （13）\(\displaystyle 2\)。 （14）\(\displaystyle \dfrac{\displaystyle 9}{\displaystyle 2}\)。


\subsection*{三、解答题}

（15）\(\displaystyle \dfrac{\displaystyle 2}{\displaystyle 3}\)。 （16）\(\displaystyle \dfrac{\displaystyle (2-\sqrt2)\pi}{\displaystyle 16}\)。 （17）\(\displaystyle \dfrac{\displaystyle 1}{\displaystyle 4}\)。 （18）\(\displaystyle (\dfrac{\displaystyle 1}{\displaystyle \ln2}-1,\allowbreak \dfrac{\displaystyle 1}{\displaystyle 2})\)。


（19）（I）证明略。（II）证明略。\(\displaystyle S(x)=\dfrac{\displaystyle e^{-x}}{\displaystyle 1-x}\)，\(\displaystyle x\in(-1,\allowbreak 1)\)。


（20）（I）证明略。（II）\(\displaystyle k(1,\allowbreak 2,\allowbreak -1)^{\mathrm T}+(1,\allowbreak 1,\allowbreak 1)^{\mathrm T}\) 为线性方程组 \(\displaystyle Ax=\beta\) 的通解，其中 \(\displaystyle k\) 为任意常数。


（21）\(\displaystyle a=2\)，\(\displaystyle Q=\begin{pmatrix}-\dfrac{\displaystyle 1}{\displaystyle \sqrt2}&\dfrac{\displaystyle 1}{\displaystyle \sqrt3}&\dfrac{\displaystyle 1}{\displaystyle \sqrt6}\\0&-\dfrac{\displaystyle 1}{\displaystyle \sqrt3}&\dfrac{\displaystyle 2}{\displaystyle \sqrt6}\\\dfrac{\displaystyle 1}{\displaystyle \sqrt2}&\dfrac{\displaystyle 1}{\displaystyle \sqrt3}&\dfrac{\displaystyle 1}{\displaystyle \sqrt6}\end{pmatrix}\)，且 \(\displaystyle Q^{\mathrm T}AQ=\begin{pmatrix}6&0&0\\0&-3&0\\0&0&0\end{pmatrix}\)。


（22）（I）\(\displaystyle \dfrac{\displaystyle 4}{\displaystyle 9}\)。（II）\(\displaystyle f_Z(z)=\begin{cases}z,\allowbreak &0<z<1,\allowbreak \\z-2,\allowbreak &2<z<3,\allowbreak \\0,\allowbreak &\text{其他}.\end{cases}\)


（23）（I）\(\displaystyle f_{Z_1}(z)=\begin{cases}\sqrt{\dfrac{\displaystyle 2}{\displaystyle \pi}}\dfrac{\displaystyle 1}{\displaystyle \sigma} e^{-z^2/(2\sigma^2)},\allowbreak &z\ge0,\allowbreak \\0,\allowbreak &z<0.\end{cases}\)（II）\(\displaystyle \hat\sigma=\sqrt{\dfrac{\displaystyle \pi}{\displaystyle 2}}\bar Z\)。（III）\(\displaystyle \hat\sigma=\sqrt{\dfrac{\displaystyle 1}{\displaystyle n}\displaystyle \sum_{i=1}^nZ_i^2}\)。


\clearpage
\end{document}
